% Source-faithful assembly of Mumford, Abelian Varieties, Chapter I,
% Sections 2--3 (source pages 27--48). Running headers and pedagogical
% asides are omitted.

\section{Line bundles: Chern classes and Appell--Humbert}
\label{mumford-source-I2-remainder}
\dcref{M:ch1:I.2}
\uses{mumford-frag-integral-cohomology}

Let $H$ be the ring of holomorphic functions on $V$ and
$H^*=H^0(V,\mathcal O_V^*)$. The exponential sequence and the quotient
description of line bundles identify the Chern class of a cocycle
$\{e_u\}\in Z^1(U,H^*)$ with the class of the integral $2$-cocycle
\[
 F(u_1,u_2)=f_{u_2}(z+u_1)-f_{u_1+u_2}(z)+f_{u_1}(z),
 \qquad e_u(z)=e^{2\pi i f_u(z)}.
\]
\source{mumford-abelian-varieties:page-0027}

\begin{lemma}[Alternation of group $2$-cocycles]
\label{mumford-source-alternation}
\dcref{M:ch1:I.2}
\uses{mumford-frag-integral-cohomology}
The map $AF(u_1,u_2)=F(u_1,u_2)-F(u_2,u_1)$ induces an isomorphism
\[
 H^2(U,\mathbf Z)\xrightarrow{\sim}\operatorname{Hom}(\Lambda^2U,\mathbf Z)
 \simeq\Lambda^2\operatorname{Hom}(U,\mathbf Z),
\]
and $A(\xi\cup\eta)=\xi\wedge\eta$.
\source{mumford-abelian-varieties:page-0027}
\end{lemma}
\begin{proof}
The cocycle identity implies, after the three cyclic substitutions and taking
the combination $(i)+(ii)-(iii)$, that
$AF(u_3,u_1+u_2)=AF(u_3,u_1)+AF(u_3,u_2)$. Since $AF$ is alternating, it is
bilinear. If $F=\delta G$, then direct cancellation gives $AF=0$, so $A$
descends to cohomology. The torus computation of integral cohomology and
compatibility of the quotient map with cup products show that the cohomology
ring of $U$ is the exterior algebra on $H^1(U,\mathbf Z)$. For homomorphisms
$f,g:U\to\mathbf Z$, the cocycle for $\xi\cup\eta$ is
$c(s,t)=f(s)g(t)$, and
\[
 A(\xi\cup\eta)(s,t)=f(s)g(t)-f(t)g(s)=(f\wedge g)(s,t).
\]
\source{mumford-abelian-varieties:page-0028}
\end{proof}

\begin{proposition}[Chern class formula]
\label{mumford-source-chern-formula}
\dcref{M:ch1:I.2}
\uses{mumford-source-alternation}
The Chern class of the line bundle defined by $\{e_u\}$ is the alternating form
\[
 E(u_1,u_2)=f_{u_2}(z+u_1)+f_{u_1}(z)-f_{u_1+u_2}(z)-f_{u_2}(z),
\]
independent of $z$, where $e_u=e^{2\pi i f_u}$.
\source{mumford-abelian-varieties:page-0029}
\end{proposition}

\begin{corollary}[Type $(1,1)$ condition]
\label{mumford-source-chern-type}
\dcref{M:ch1:I.2}
\uses{mumford-source-chern-formula}
After extending $E$ $\mathbf R$-linearly to $V$, one has
$E(ix,iy)=E(x,y)$.
\source{mumford-abelian-varieties:page-0029}
\end{corollary}
\begin{proof}
The image of a Chern class in $H^2(X,\mathcal O_X)$ vanishes. Under
\[
 H^2(X,\mathbf C)\simeq(\Lambda^2T)\oplus(T\otimes\overline T)
 \oplus(\Lambda^2\overline T),
\]
the map to $H^2(X,\mathcal O_X)\simeq\Lambda^2\overline T$ is projection.
Reality gives $E_1=\overline{E_2}$, so vanishing is equivalent to $E$ lying
in $T\otimes\overline T$, which is equivalent to $E(ix,iy)=E(x,y)$.
\end{proof}


\begin{lemma}[Hermitian forms and invariant alternating forms]
\label{mumford-source-hermitian-correspondence}
\dcref{M:ch1:I.2}
\uses{mumford-source-chern-type}
Hermitian forms $H$ on $V$ correspond bijectively to real skew-symmetric forms
$E$ satisfying $E(ix,iy)=E(x,y)$ by
\[
 E(x,y)=\operatorname{Im}H(x,y),\qquad H(x,y)=E(ix,y)+iE(x,y).
\]
\source{mumford-abelian-varieties:page-0030}
% Source explicitly leaves this elementary correspondence proof to the reader.
\notready
\end{lemma}

\begin{lemma}[Semicharacters]
\label{mumford-source-semicharacter}
\dcref{M:ch1:I.2}
\uses{mumford-source-hermitian-correspondence}
If $H$ is Hermitian, $E=\operatorname{Im}H$ is integral on $U$, then there is
$\alpha:U\to\mathbf C_1^*$ with
\[
 \alpha(u_1+u_2)=e^{i\pi E(u_1,u_2)}\alpha(u_1)\alpha(u_2),
\]
and
\[
 e_u(z)=\alpha(u)e^{\pi H(z,u)+\frac12\pi H(u,u)}
\]
is a $1$-cocycle whose Chern class is $E$.
\source{mumford-abelian-varieties:page-0031}
\end{lemma}
\begin{proof}
Writing the linear solutions as $f_u(z)=H(z,u)/(2i)+\beta_u$ and imposing
the cocycle equation gives, after a linear coboundary normalization, pure
imaginary constants $\gamma_u$. Setting $\alpha(u)=e^{2\pi\gamma_u}$ gives
the displayed semicharacter equation and $|\alpha(u)|=1$; substitution gives
the displayed $e_u$ and the cocycle identity.
\end{proof}

\begin{definition}[The bundle $L(H,\alpha)$]
\label{mumford-source-LHa}
\dcref{M:ch1:I.2}
\uses{mumford-source-semicharacter}
$L(H,\alpha)$ is the quotient of $\mathbf C\times V$ by
\[
 (\lambda,z)\longmapsto
 (\alpha(u)e^{\pi H(z,u)+\frac12\pi H(u,u)}\lambda,z+u).
\]
Moreover $L(H_1,\alpha_1)\otimes L(H_2,\alpha_2)\simeq
L(H_1+H_2,\alpha_1\alpha_2)$.
\source{mumford-abelian-varieties:page-0031}
\end{definition}

\begin{theorem}[Appell--Humbert]
\label{mumford-thm-appell-humbert}
\dcref{M:ch1:I.2}
Every line bundle on $X=V/U$ is uniquely $L(H,\alpha)$ as above. There are
exact sequences
\[
0\to\operatorname{Hom}(U,\mathbf C^*)\to\{(H,\alpha)\}\to
\{H:(\operatorname{Im}H)(U\times U)\subset\mathbf Z\}\to0,
\]
and
\[
0\to\Pic^0X\to\Pic X\xrightarrow{c^1}
\ker[H^2(X,\mathbf Z)\to H^2(X,\mathcal O_X)]\to H^2(X,\mathcal O_X)\to0.
\]
\source{mumford-abelian-varieties:page-0031}\source{mumford-abelian-varieties:page-0032}
\end{theorem}
\begin{proof}
The type $(1,1)$ corollary identifies the image of the Chern map with the
Hermitian forms. If $\alpha\in\operatorname{Hom}(U,\mathbf C_1^*)$ becomes a
coboundary, $g(z+u)/g(z)=\alpha(u)$. A compact fundamental set and
$|\alpha|=1$ force $g$ constant, so $\alpha=1$. The exponential-sequence
diagram, together with the surjectivity of $H^1(X,\mathbf C)\to
H^1(X,\mathcal O_X)$, proves surjectivity. A coboundary normalization makes
every character unitary, completing the classification.
\end{proof}

\section{Appendix: group cohomology}
\label{mumford-source-group-cohomology}
\dcref{M:ch1:I.2}

For a $G$-module $M$, let $C^p(G,M)$ be functions $G^p\to M$, with
\[
\delta f(\sigma_0,\ldots,\sigma_p)=\sigma_0f(\sigma_1,\ldots,\sigma_p)
 +\sum_{i=0}^{p-1}(-1)^{i+1}f(\sigma_0,\ldots,\sigma_i\sigma_{i+1},\ldots,\sigma_p)
 +(-1)^{p+1}f(\sigma_0,\ldots,\sigma_{p-1}).
\]
Write $H^p(G,M)=\ker\delta/\operatorname{im}\delta$. Cup products are given
by
\[
 (f\cup g)(\sigma_1,\ldots,\sigma_{p+q})=
 f(\sigma_1,\ldots,\sigma_p)*\sigma_1\cdots\sigma_p
 g(\sigma_{p+1},\ldots,\sigma_{p+q}).
\]
\source{mumford-abelian-varieties:page-0033}

\begin{theorem}[Quotient-space cohomology map]
\label{mumford-source-quotient-cohomology}
\dcref{M:ch1:I.2}
For $Y=X/G$ with $G$ free and discontinuous and $\pi:X\to Y$, every sheaf
$\mathcal F$ on $Y$ has a natural map
\[
 H^p(G,\Gamma(X,\pi^*\mathcal F))\to H^p(Y,\mathcal F).
\]
It is compatible with exact sequences and cup products, and is an isomorphism
if $H^i(X,\pi^*\mathcal F)=0$ for $i\ge1$.
\source{mumford-abelian-varieties:page-0033}\source{mumford-abelian-varieties:page-0034}\source{mumford-abelian-varieties:page-0035}
\end{theorem}
\begin{proof}
Choose a covering $V_i$ with lifts $U_i$ and unique transition elements
$\sigma_{ij}$. Map group cochains to Cech cochains by
\[
 (\phi_pf)_{i_0\ldots i_p}=\operatorname{res}(\pi^*)_{i_0}^{-1}
 f(\sigma_{i_0i_1},\ldots,\sigma_{i_{p-1}i_p}).
\]
Direct calculation gives $\delta\phi_p=\phi_{p+1}\delta$ and proves the
compatibilities. For the vanishing assertion, embed $\mathcal F$ in an
injective $\mathcal F'$, set $\mathcal F''=\mathcal F'/\mathcal F$, and use
induction on $p$ in the resulting long exact sequences; exactness of global
sections follows from $H^1(X,\pi^*\mathcal F)=0$.
\end{proof}

\section{Algebraizability of tori}
\label{mumford-source-I3}
\dcref{M:ch1:I.3}
\uses{mumford-source-LHa}

For $L(H,\alpha)$, sections correspond to holomorphic functions $\theta$ on
$V$ satisfying
\[
 \theta(z+u)=\alpha(u)e^{\pi H(z,u)+\frac12\pi H(u,u)}\theta(z),
 \qquad u\in U.
\]
Such a function is called a theta-function for $(H,\alpha)$.
\source{mumford-abelian-varieties:page-0036}

\begin{proposition}[Degenerate and negative Hermitian forms]
\label{mumford-source-theta-degenerate}
\dcref{M:ch1:I.3}
\uses{mumford-source-LHa}
Let $N=\{x:H(x,y)=0\ \forall y\}$.  If $H$ is degenerate, every
theta-function is constant on cosets of $N$, and a nonzero theta-function can
exist only if $\alpha|_{N\cap U}=1$; the associated sections factor through
$X/(N/(N\cap U))$, so $L(H,\alpha)$ is not ample.  If $H$ is negative on a
nonzero complex subspace, then it has no nonzero theta-functions.
\source{mumford-abelian-varieties:page-0036}\source{mumford-abelian-varieties:page-0037}
\end{proposition}
\begin{proof}
The identities $N=\{x:E(x,y)=0\ \forall y\}$ and integrality of $E$ show
that $N\cap U$ is a lattice in $N$.  For $u\in N\cap U$ the functional
equation gives $\theta(z+u)=\alpha(u)\theta(z)$.  If $K\subset N$ is compact
with $N=K+(N\cap U)$, then
\[
 |\theta(z_0+z')|\leq\sup_{\zeta\in K}|\theta(z_0+\zeta)|.
\]
The maximum principle makes $\theta$ constant on every $N$-coset, proving
the factorization and the ampleness assertion.

Now suppose $H(w,w)<0$ on a complex subspace $W$.  Choose compact $K$ with
$V=U+K$, write $w=d+u$ with $d\in K$, and use the functional equation:
\[
 |\theta(z_0+w)|=|\theta(z_0+d)|
 e^{\pi\operatorname{Re}H(z_0+d,w)+\frac12\pi H(w,w)}.
\]
The exponent equals
\[
 \tfrac12H(w,w)+\operatorname{Re}H(z_0,w)+c(d,z_0),
\]
where $c(d,z_0)$ is bounded on $K$.  It tends to $-\infty$ along $W$, so
the maximum principle applied on $W$ gives $\theta(z_0+w)=0$ for every $w$;
hence $\theta=0$.
\end{proof}

\begin{proposition}[Dimension of theta-functions]
\label{mumford-source-theta-dimension}
\dcref{M:ch1:I.3}
\uses{mumford-source-LHa}
If $H$ is positive definite and $E=\operatorname{Im}H$ is integral on $U$,
then
\[
 \dim H^0(X,L(H,\alpha))=\sqrt{\det_U E}.
\]
\source{mumford-abelian-varieties:page-0037}\source{mumford-abelian-varieties:page-0038}\source{mumford-abelian-varieties:page-0039}\source{mumford-abelian-varieties:page-0040}
\end{proposition}
\begin{proof}
Choose a rank-$g$ sublattice $U'\subset U$ with $E(U',U')=0$ and, for
$W=\mathbf R\cdot U'$, $W\cap iW=0$.  Let $B$ be the symmetric complex
bilinear form extending $H|_{W\times W}$, and choose a complex-linear
$\lambda$ real on $W$ with $\alpha(u)=e^{2\pi i\lambda(u)}$ on $U'$.  Then
$e^{-2\pi i\lambda(z)}e^{-\frac12\pi B(z,z)}\theta(z)$ is $U'$-periodic,
so it has a Fourier expansion
\[
 \theta^*(z)=\sum_{\chi\in\widehat U'}c_\chi
 e^{2\pi i(\chi(z)+\lambda(z))}. \tag{1}
\]
For $u\in U$, define $\widehat u\in\widehat U'$ by
$\widehat u(u')=E(u',u)$.  Since
$(H-B)(z,u)=2i\widehat u(z)$, substitution of (1) in the functional
equation gives
\[
 c_\chi=\alpha(u)e^{i\pi\widehat u(u)-2\pi i[\chi(u)+\lambda(u)]}
 c_{\chi-\widehat u}. \tag{2}
\]
Thus coefficients are determined by representatives of
$\widehat U'/M$, where $M$ is the image of $U\to\widehat U'$, and every
choice of representative coefficients extends uniquely.

For completeness, take one representative $\chi_0$ and set all other
representative coefficients to zero.  On a compact $K\subset V$, the
corresponding series is bounded by
\[
 \operatorname{const.}\sum_{\widehat u\in M}|c_{\chi_0+\widehat u}|
 e^{2\pi|\widehat u(z)|}
 \leq\operatorname{const.}\sum_{\widehat u\in M}
 e^{\pi\operatorname{Im}\widehat u(u)+A|\widehat u|^2}.
\]
Writing $z=\phi(z)+i\psi(z)$ with $\phi,\psi:V\to W$ real-linear gives
\[
 \operatorname{Im}\widehat u(u)=-H(\psi(u),\psi(u)),
\]
a negative definite quadratic form on $M$.  The series therefore converges
normally, and the dimension is $[\widehat U':M]$.

Finally, if $n$ is the order of the cokernel of
$U\to\operatorname{Hom}_{\mathbf Z}(\widehat U',\mathbf Z)$, the exact
diagram
\[
\begin{array}{ccccccccc}
0&\to&U'&\to&U&\to&U/U'&\to&0\\
&&\downarrow&&\downarrow&&\downarrow\\
0&\to&(U'/U')^\wedge&\to&\widehat U'&\to&\widehat{U'}&\to&0
\end{array}
\]
shows that the first and third vertical maps have cokernels of the same order,
whereas the middle cokernel has order $|\det E|$.  Hence
$|\det E|=n^2$ and $[\widehat U':M]=\sqrt{\det E}$.
\end{proof}

\begin{theorem}[Lefschetz]
\label{mumford-source-lefschetz}
\dcref{M:ch1:I.3}
\uses{mumford-source-theta-degenerate,mumford-source-theta-dimension}
For $L=L(H,\alpha)$ on $X=V/U$, the following are equivalent:
\begin{enumerate}
\item for every complex subtorus $Y$ there are $N>0$, a section $\sigma$ of
  $L^N$, and $x_1,x_2$ with $x_1-x_2\in Y$, $\sigma(x_1)=0$ and
  $\sigma(x_2)\ne0$;
\item $H$ is positive definite;
\item the sections of $L^n$ embed $X$ as a closed complex submanifold of a
  projective space for every $n\ge3$.
\end{enumerate}
\source{mumford-abelian-varieties:page-0040}\source{mumford-abelian-varieties:page-0041}\source{mumford-abelian-varieties:page-0042}\source{mumford-abelian-varieties:page-0043}\source{mumford-abelian-varieties:page-0044}
\end{theorem}
\begin{proof}
We have already proved (1)$\Rightarrow$(2), and (3)$\Rightarrow$(1) is
clear.  Assume $H$ positive definite and prove (3) for $n=3$; the same
argument applies for $n>3$.  If $\theta$ is a section of $L(H,\alpha)$ and
$a,b\in V$, then
\[
 \theta(z-a)\theta(z-b)\theta(z+a+b)
\]
is a section of $L^3$, since its factor of automorphy is
$\alpha(u)^3e^{3\pi H(z,u)+\frac32\pi H(u,u)}$.  The dimension proposition
gives a nonzero $\theta$, and $a,b$ can be chosen so this product does not
vanish at any prescribed point.  Thus a basis of $H^0(X,L^3)$ defines a
holomorphic map $\Theta:X\to\mathbf P^d$.

Suppose $\Theta(z_1)=\Theta(z_2)$ with $\sigma=z_2-z_1\notin U$.  For every
theta-function $\theta$ and all $a,b$, comparison of the two products and
logarithmic differentiation in $a$ gives
\[
 -\omega(z_1-a)+\omega(z_1+a+b)
 =-\omega(z_2-a)+\omega(z_2+a+b),
\]
where $\omega=d\theta/\theta$.  Hence
$\omega(z_2+z)-\omega(z_1+z)=dl(z)$ for a complex-linear form $l$, and
\[
 \theta(z+\sigma)=Ae^{l(z)}\theta(z).
\]
Comparing this with the functional equation for $u\in U$ gives
$e^{\pi H(\sigma,u)}=e^{l(u)}$, hence $\pi H(u,\sigma)=l(u)$ and
$E(\sigma,u)\in\mathbf Z$.  Therefore $\sigma\in U^+$, where
$U^+=\{x:E(x,u)\in\mathbf Z\ \forall u\in U\}$, and $\theta$ is a
theta-function for the larger lattice $U'=U+\mathbf Z\sigma$.  But the
dimension formula gives $\sqrt{\det_{U'}E}>\sqrt{\det_UE}$, while only
finitely many multipliers on $U'$ extend $\alpha$; not every $\theta$ can
extend, a contradiction.  Thus $\Theta$ is injective.

If $d\Theta$ has a nonzero kernel at $z_0$, there are $\alpha_0,\alpha_i$,
not all $\alpha_i$ zero, such that for every section $\phi$ of $L^3$,
\[
 \alpha_0\phi(z_0)+\sum_i\alpha_i\frac{\partial\phi}{\partial z_i}(z_0)=0.
\]
With $D=\sum_i\alpha_i\partial/\partial z_i$ and the triple products above,
$f=D(\log\theta)$ satisfies
$f(z_0-a)+f(z_0-b)+f(z_0+a+b)=-\alpha_0$.  Hence $f$ is affine-linear, and
integration gives a nonzero $\alpha\in V$ with
\[
 \theta(z+\lambda\alpha)=e^{c\alpha^2+\lambda f(z)}\theta(z).
\]
Comparing the transformation laws again shows that every $\lambda\alpha$
lies in the discrete lattice $U^+$, impossible for all $\lambda\in\mathbf C$.
Thus $\Theta$ is an immersion.  Since $X$ is compact, an injective immersion
is a closed embedding.
\end{proof}

\begin{theorem}[Chow]
\label{mumford-source-chow}
\dcref{M:ch1:I.3}
\uses{mumford-thm-appell-humbert}
If $X$ is complete algebraic and $Y$ is a closed analytic subset of
$X_{\mathrm{hol}}$, then $Y$ is Zariski closed. A holomorphic map between
complete algebraic varieties is algebraic, and meromorphic functions on a
complete algebraic variety are rational.
\source{mumford-abelian-varieties:page-0044}\source{mumford-abelian-varieties:page-0045}
\end{theorem}
\begin{proof}
Chow's theorem for projective space, together with Chow's lemma, proves the
closed-analytic-subset assertion.  If $f:X_{\mathrm{hol}}\to Y_{\mathrm{hol}}$
is holomorphic, its graph is a closed analytic subset of
$(X\times Y)_{\mathrm{hol}}$, hence a closed algebraic subset of $X\times Y$.
At $(x,f(x))$, write $\mathcal O_1=\mathcal O_{X,x}$,
$\mathcal O_2=\mathcal O_{Y,f(x)}$ and use the analytic local rings
$\widetilde{\mathcal O}_i$.  The graph projection gives
\[
\begin{array}{ccc}
\mathcal O_1&\longrightarrow&\mathcal O_2\\
\downarrow&&\downarrow\\
\widetilde{\mathcal O}_1&\xrightarrow{\sim}&\widetilde{\mathcal O}_2.
\end{array}
\]
Reducing modulo squares and applying Nakayama first yields
$\mathfrak m_2=\mathfrak m_1\mathcal O_2$, then
$\mathcal O_2=\mathcal O_1$.  Hence the graph projection is an algebraic
isomorphism and $f$ is algebraic.  The same graph argument applied to a
meromorphic function shows it lies in $\mathbf C(X)$.
\end{proof}

\begin{corollary}[Algebraizability criterion]
\label{mumford-source-algebraizability}
\dcref{M:ch1:I.3}
For a complex torus $X=V/U$, the following are equivalent: $X$ is associated to
a projective algebraic variety; $X$ is associated to an algebraic variety;
$X$ has $g$ algebraically independent meromorphic functions; and there is a
positive definite Hermitian $H$ with $\operatorname{Im}H$ integral on $U$.
\source{mumford-abelian-varieties:page-0045}\source{mumford-abelian-varieties:page-0046}
\end{corollary}
\begin{proof}
The implications (1)$\Rightarrow$(2)$\Rightarrow$(3) are immediate, and
(4)$\Rightarrow$(1) is Lefschetz.  For (3)$\Rightarrow$(4), let
$f_1,\ldots,f_g$ be algebraically independent meromorphic functions, let
$D_i$ be their polar divisors, $D=\sum D_i$, and let $L=\mathcal O_X(D)$.
There are sections $\sigma_0,\ldots,\sigma_g$ with
$\sigma_i=f_i\sigma_0$ where $f_i$ is regular.  Appell--Humbert gives
$L=L(H,\alpha)$ with $H$ semidefinite.  Let $V_0$ be its nullspace and
$X_0=V_0/(U\cap V_0)$.  The induced form on $X/X_0$ is positive definite,
and zeros of sections of $L$ are unions of $X_0$-cosets; hence each $f_i$
descends to an algebraically independent $\overline f_i$ on $X/X_0$.  Thus
\[
 g\le\operatorname{tr.deg}_{\mathbf C}\mathbf C(X/X_0)
 =\dim(X/X_0)\le g,
\]
so $X_0=0$ and $H$ is positive definite.
\end{proof}

\begin{lemma}[Generic nonalgebraizability]
\label{mumford-source-generic-nonalgebraic}
\dcref{M:ch1:I.3}
For almost all lattices $U$ in a $g$-dimensional complex vector space, the map
\[
 \Lambda^2_{\mathbf Z}U\longrightarrow\Lambda^2_{\mathbf C}V
\]
is injective. Hence almost all complex tori of dimension at least $2$ are
nonalgebraic.
\source{mumford-abelian-varieties:page-0047}\source{mumford-abelian-varieties:page-0048}
% Source explicitly leaves this genericity lemma proof to the reader.
\notready
\end{lemma}
